\section{Preisfindung auf dem Strommarkt}

\begin{frame}{Grundsätzliches}
    \mytext{
    \begin{itemize}
        \item Liberalisierte Strommärkte: 
        \item[] jeder darf verkaufen und kaufen
        \item OTC Over The Counter
        \item börslicher Handel: EEX/EPEX
        \item Spotmärkte, Terminmärkte, Derivate
        \item Händler als eigene Marktrolle (Dienstleister)
        \item Bilanzkreisverantwortliche Händler
        \item Zertifikatehandel (\ce{CO2} etc.)
        \item Verbindung zu internationalen Energiemärkten (fossil)
    \end{itemize}
    }
\end{frame}

% \begin{frame}{Ökonomische Grundlagen}

%     {\color{red}
%     (Umfang und Vorwissen klären, wird noch eingestellt)
%     }
% \end{frame}

\begin{frame}{Merit Order Effekt}
    \mycomment{
    Kraftwerksbetreiber akzeptieren jeden Preis, der oberhalb ihrer Grenzkosten liegt. Diese sind bei Erneuerbaren besonders niedrig. Ordnet man diese Gebote von links nach rechts, so bekommt man die \enquote{Merit Order} genannte Reihenfolge.

    In dieser Darstellung übersteigt die Nachfrage das Angebot von Strom durch Erneuerbare. Verbraucher sind bereit, mehr zu zahlen, nämlich den Preis des Stromangebots durch Kohlekraftwerke. 
    }
    \showmytikz[1]{../shared/tikz/tikz-meritorder-de}
\end{frame}

\begin{frame}{Merit Order Effekt}
    \mycomment{
    Da dieser Preis höher liegt als die Grenzkosten der Erneuerbaren, machen die Erneuerbaren gute Gewinne. Dies ist eine Form der Knappheitsrente.
    }
    \showmytikz[1,2]{../shared/tikz/tikz-meritorder-de}
\end{frame}

\begin{frame}{Merit Order Effekt}
    \mycomment{
    Steigt das Angebot an günstiger Energie, so fällt irgendwann der Preis. 
    
    Da die Einspeisung aus Erneuerbaren stark von der Saison, Tageszeit und Wetterbedingungen abhängt, schwankt die von ihnen angebotene Menge.
    
    Da die Angebotskurve nicht stetig ist, sondern die Merit Order einer Treppe gleicht, kommt es zu starken Preisschwankungen bzw. -sprüngen.

    }
    \showmytikz[3]{../shared/tikz/tikz-meritorder-de}
\end{frame}
